\chapter{Metodologie výzkumu}
\label{chap:methodology}

Tato kapitola popisuje výzkumný design, výběr metod sběru dat, způsob výběru respondentů a~etické aspekty realizovaného výzkumu. Zvolený přístup vychází z~principů smíšené metodologie, která kombinuje kvalitativní a~kvantitativní výzkumné metody za účelem komplexního pochopení zkoumaného problému \parencite{creswell_2018}.

% ============================================================================
\section{Výzkumný design}
\label{sec:meth_design}

Pro zodpovězení stanovených výzkumných otázek byl zvolen sekvenční výzkumný design, ve kterém kvalitativní fáze předchází fázi kvantitativní. Tento přístup umožňuje nejprve prostřednictvím skupinových diskusí získat hlubší vhled do motivací a~bariér spotřebitelů a~následně ověřit získané poznatky na reprezentativnějším vzorku prostřednictvím dotazníkového šetření.

Výzkum se skládá ze tří hlavních fází:
\begin{enumerate}
  \item \textbf{Kvalitativní fáze} — skupinové diskuse (focus groups),
  \item \textbf{Senzorická analýza} — Central Location Test (\gls{clt}),
  \item \textbf{Kvantitativní fáze} — online dotazníkové šetření (\gls{cawi}).
\end{enumerate}

% ============================================================================
\section{Kvalitativní fáze — Focus Groups}
\label{sec:meth_focus_groups}

Skupinové diskuse (focus groups) představují jednu z~nejrozšířenějších kvalitativních metod marketingového výzkumu. \textcite{krueger_casey_2015} je definují jako pečlivě plánované diskuse v~přátelském prostředí, jejichž cílem je získat vnímání cílové skupiny na definované téma.

\subsection{Parametry skupinových diskusí}

Budou realizovány dvě skupinové diskuse, každá o~7--9 respondentech, v~délce přibližně 80 minut:

\begin{table}[H]
\centering
\caption{Parametry skupinových diskusí}
\label{tab:fgd_params}
\begin{tabular}{lll}
\toprule
\textbf{Parametr} & \textbf{Skupina 1} & \textbf{Skupina 2} \\
\midrule
Cílová skupina & Gen Z a mileniálové & Matky s dětmi \\
Věk & 18--40 let & 25--40 let \\
Počet respondentů & 7--9 & 7--9 \\
Délka & 80 minut & 80 minut \\
\bottomrule
\end{tabular}
\end{table}

\subsection{Struktura diskuse}

Diskuse je rozdělena do tří tematických bloků:

\begin{description}
  \item[Blok I — Vnímání produktu a~první dojem:] Mapování spontánních asociací s~kategorií mraženého ovoce v~čokoládě, povědomí o~značce Franui, emocionální reakce na koncept produktu.
  \item[Blok II — Senzorické testování:] Řízená ochutnávka, hodnocení textury, chuťové rovnováhy a~senzorického kontrastu v~reálném čase.
  \item[Blok III — Design a~positioning:] Hodnocení vizuální stránky obalu, vnímání ceny, důvěra v~lokální původ a~\gls{bio} kvalitu, zájem o~budoucí varianty.
\end{description}

% ============================================================================
\section{Senzorická analýza — Central Location Test}
\label{sec:meth_clt}

Senzorická analýza bude realizována formou \gls{clt}, což je náročný typ výzkumu vyžadující vysokou míru soustředění respondenta \parencite{lawless_heymann_2010}. Testování proběhne za kontrolovaných podmínek s~využitím zaškolených tazatelů.

\subsection{Hodnocené parametry}

Senzorická analýza se zaměří na tři klíčové oblasti:

\begin{enumerate}
  \item \textbf{Hodnocení textury v~čase:} Sledování změn struktury ovoce a~čokolády v~intervalech 5, 10 a~15 minut od vytažení z~mrazáku. Tímto způsobem se ověří kvalita technologie šokového mražení, která je klíčová pro zachování struktury ovoce \parencite{grover_negi_2023}.
  \item \textbf{Chuťová rovnováha a~kontrast:} Posuzování poměru mezi sladkostí čokolády a~kyselostí ovoce, hodnocení celkového senzorického dojmu.
  \item \textbf{Otevřené hodnocení:} Volné popisy senzorických vjemů za asistence zaškolených tazatelů.
\end{enumerate}

% ============================================================================
\section{Kvantitativní fáze — CAWI}
\label{sec:meth_cawi}

Kvantitativní fáze výzkumu bude realizována prostřednictvím online dotazníkového šetření metodou \gls{cawi}. Tato metoda byla zvolena pro svou efektivitu při oslovení širokého spektra respondentů a~rychlost sběru dat \parencite{punch_2014}.

\subsection{Parametry dotazníkového šetření}

\begin{table}[H]
\centering
\caption{Parametry kvantitativního výzkumu}
\label{tab:cawi_params}
\begin{tabular}{ll}
\toprule
\textbf{Parametr} & \textbf{Hodnota} \\
\midrule
Metoda sběru dat & CAWI (online dotazník) \\
Cílová skupina & 18--45 let \\
Výběrové kritérium & Aktivní uživatelé Instagramu a TikToku \\
Velikost vzorku & 400 respondentů \\
Reprezentativita & Kvótní výběr (věk, pohlaví, region) \\
\bottomrule
\end{tabular}
\end{table}

\subsection{Tematické okruhy dotazníku}

Dotazník je strukturován do tří hlavních tematických okruhů:

\begin{enumerate}
  \item \textbf{Chuťové preference:} Zjišťování preferencí ohledně druhů ovoce, typů čokolády a~jejich kombinací.
  \item \textbf{Nákupní chování:} Příležitosti a~motivace k~nákupu mraženého ovoce v~čokoládě, distribuční preference.
  \item \textbf{Cenová citlivost:} Aplikace metody \gls{psm} \parencite{vanwestendorp_1976} pro stanovení optimální cenové hladiny.
\end{enumerate}

\subsection{Metoda Price Sensitivity Meter}

Pro měření cenové citlivosti bude využita metoda \gls{psm}. Respondenti budou odpovídat na čtyři standardizované otázky (viz Kapitola~\ref{sec:theory_pricing}), jejichž výsledky umožní identifikovat:
\begin{itemize}
  \item optimální cenový bod (OPP — Optimal Price Point),
  \item bod cenové indiference (IDP — Indifference Price Point),
  \item přijatelné cenové rozmezí.
\end{itemize}

% ============================================================================
\section{Výběr respondentů}
\label{sec:meth_sampling}

Pro kvalitativní fázi budou respondenti vybíráni účelovým výběrem (purposive sampling) s~ohledem na definované segmenty — generace~Z/mileniálové a~matky s~dětmi \parencite{morgan_1997}.

Pro kvantitativní fázi bude uplatněn kvótní výběr s~cílem zajistit, aby struktura vzorku odpovídala reálnému rozložení české populace ve věku 18--45 let podle věku, pohlaví a~regionálního rozložení \parencite{czso_spotreba_2024}.

% ============================================================================
\section{Etické aspekty výzkumu}
\label{sec:meth_ethics}

Výzkum bude realizován v~souladu s~etickými principy výzkumné práce:
\begin{itemize}
  \item respondenti budou informováni o~účelu výzkumu a~způsobu využití dat,
  \item účast ve výzkumu bude dobrovolná,
  \item osobní údaje budou anonymizovány,
  \item respondenti v~rámci senzorického testování budou předem informováni o~složení testovaných produktů s~ohledem na případné potravinové alergie.
\end{itemize}
